\begin{proposition}\label{prop:frontier-polyhedral-bergman-positive-infinity-obstruction}
In the notation of \cite[Example~6.24 and Claim~6.25]{BMPS26}, let
\(g\geq\gamma\) be its obstacle and put \(u=g-\gamma\).  Write
\(\ell(x_1,x_2,x_3)=x_1+x_2\), so that
\[
 \ell(e_{12})=2,
 \qquad
 \ell(e_{34})=-2,
 \qquad
 \bigl(\ell(e_1),\ell(e_2),\ell(e_3),\ell(e_4)\bigr)
   =(1,1,0,-2).
\]
For \(\lambda\in\mathbf Q\cap(0,\tfrac12)\) and
\(R\in\mathbf Q\) with \(R>2\lambda\varepsilon\), set
\[
 \gamma_{\lambda,R}
   :=\max\{\gamma,\lambda\ell-R,-\lambda\ell-R\}.
\]
Then \(\gamma_{\lambda,R}\) is piecewise-affine polyhedrally convex and
positive at infinity.  If
\[
 P:=P_{\gamma_{\lambda,R}}(u),
\]
then \(P=\gamma\) on a neighbourhood of the origin.  Moreover, either
\[
 P\notin
 \mathrm{PC}_{\mathrm{reg},\mathbf R}
     (\mathbf X,\gamma_{\lambda,R}),
\]
or
\[
 \int_{\mathbf X}
 \bigl(P-(\gamma_{\lambda,R}+u)\bigr)\,
 \MA_{\mathrm{poly},\mathbf R}(P)
 \leq -4\varepsilon<0.
\]
Consequently \((\mathbf X,\gamma_{\lambda,R})\) cannot have both
real-polyhedral regularity and real-polyhedral orthogonality.  The same
dichotomy precludes simultaneous literal regularity and orthogonality.
\end{proposition}

\begin{proof}
The restrictions of the two affine functions in the displayed maximum are
polyhedrally convex, as is \(\gamma\); a finite maximum preserves the
polyhedral-convex-combination inequalities.  Thus
\(\gamma_{\lambda,R}\) is piecewise-affine polyhedrally convex by
\cite[Definitions~3.1 and~3.3 and Lemma~3.4]{BMPS26}.  On every recession
direction \(v\), its recession slope is
\[
 \max\{\gamma^{\infty}(v),\lambda|\ell(v)|\}.
\]
The source's ray values show that \(\gamma^{\infty}\) can vanish only on
the line \(\mathbf R e_{12}=\mathbf R e_{34}\), where
\(|\ell|\) is nonzero.  This proves positivity at infinity.

Set
\[
 \gamma_{\infty}:=\max\{\gamma,\lambda\ell,-\lambda\ell\}.
\]
This is a conical piecewise-affine polyhedrally convex function on the
conical fan, and
\[
 \gamma_{\infty}-R\leq\gamma_{\lambda,R}\leq\gamma_{\infty}.
 \tag{\(\dagger\)}\label{eq:frontier-polyhedral-bergman-bounded-reference-shift}
\]
Thus \(\delta:=\gamma_{\lambda,R}-\gamma_{\infty}\) lies in
\(\mathrm{PA}_{b}(\mathbf X)\).  For every
\(v\in\mathrm{PA}_{b}(\mathbf X)\),
\Cref{prop:frontier-polyhedral-bounded-reference-shift} gives
\[
 P_{\gamma_{\lambda,R}}(v)
   =P_{\gamma_{\infty}}(v+\delta).
 \tag{\(\ddagger\)}\label{eq:frontier-polyhedral-bergman-envelope-shift}
\]
This identity is only a reference-change formula.  Remark~6.6 of the source
regularizes conical functions; it does not prove regularity for the
nonconical obstacle \(v+\delta\).

Since \(u\geq0\), \(\gamma_{\lambda,R}\) is a competitor for its own
envelope, hence
\[
 \gamma_{\lambda,R}\leq P\leq\gamma_{\lambda,R}+u.
 \tag{\(*\)}\label{eq:frontier-polyhedral-bergman-growth-envelope-bounds}
\]
On both exceptional rays beyond the source's length-\(\varepsilon\)
modification, \(u=0\).  Choose
\(a\) with \(\varepsilon<a<R/(2\lambda)\).  At the opposite points
\(ae_{12}\) and \(-ae_{12}=ae_{34}\), the two bounds in
\eqref{eq:frontier-polyhedral-bergman-growth-envelope-bounds} agree and
are zero.  Convexity on their connecting line, together with the lower
bound, therefore gives
\[
 P=\gamma=0
 \quad\text{on }[-ae_{12},ae_{12}],
 \qquad\text{in particular }P(0)=0.
 \tag{\(**\)}\label{eq:frontier-polyhedral-bergman-central-contact}
\]
On every other primitive ray of the Bergman fan, the choice
\(\lambda<\tfrac12\) gives
\(\gamma_{\lambda,R}=\gamma\), while the source obstacle has
\(u=\varepsilon\).  The restriction of \(P\) to such a ray is convex,
starts at zero by \eqref{eq:frontier-polyhedral-bergman-central-contact},
and lies between \(t\) and \(t+\varepsilon\).  Its excess over \(t\)
cannot be positive at any point: a positive secant slope would persist by
convexity and contradict boundedness by \(\varepsilon\).  Hence it equals
\(\gamma\) on every such ray.  On the two exceptional rays it equals
\(\gamma\) on the central segment by
\eqref{eq:frontier-polyhedral-bergman-central-contact}.

Near the origin the two affine terms \(\pm\lambda\ell-R\) are negative,
so \(\gamma_{\lambda,R}=\gamma\).  In each of the six two-dimensional
cones, a point sufficiently near the origin lies on a segment joining a
boundary ray to the added diagonal ray.  The source reference \(\gamma\)
is affine on the corresponding side of that diagonal.  Applying convexity
of \(P\) to this segment and using the ray equalities just proved gives
\(P\leq\gamma\); the lower bound in
\eqref{eq:frontier-polyhedral-bergman-growth-envelope-bounds} gives the
reverse inequality.  Thus \(P=\gamma\) on a neighbourhood of \(0\).

If
\(P\notin\mathrm{PC}_{\mathrm{reg},\mathbf R}
(\mathbf X,\gamma_{\lambda,R})\), then this obstacle already disproves
real-polyhedral regularity.  Otherwise its real-scalar Monge--Amp\`ere
measure is defined.  Locality and the source computation
\(\MA_{\mathrm{poly},\mathbf R}(\gamma)\{0\}=4\) give
\(\MA_{\mathrm{poly},\mathbf R}(P)\{0\}=4\).  Finally \(P\leq
\gamma_{\lambda,R}+u\) everywhere, whereas
\(P(0)=0<\gamma_{\lambda,R}(0)+u(0)=\varepsilon\).  The displayed
integral is therefore at most \(-4\varepsilon\), which is incompatible
with real-polyhedral orthogonality.  Since the literal regularizable class
is contained in the real one and the measures agree on literal
approximants, the same alternative rules out the simultaneous literal
properties.
\end{proof}